\section{Sikkerhetskrav}
\label{sec:krav-sikkerhet}

VideoAPI opererer i et kritisk helsedomene hvor uautorisert tilgang til 
videostrømmer kan utgjøre alvorlige personvernbrudd. Sikkerhetskravene 
sikrer at kun autoriserte brukere kan opprette, administrere og delta i 
videorom.

\subsection{API-autentisering (S1.1)}
\label{subsec:krav-s11}

API-et skal være beskyttet med token-basert autentisering (OAuth 2.0), 
der alle endepunkter (unntatt helsesjekk-endepunkter) krever gyldig 
tilgangstoken.

\textbf{Rasjonale:} Kun autoriserte AMK/LVS-systemer skal kunne opprette 
og administrere videorom. Token-basert autentisering sikrer at kun HDOs 
egne systemer kan kommunisere med VideoAPI.

\textbf{Implementasjon:} VideoAPI bruker JWT Bearer Authentication for 
å validere innkommende forespørsler.

\subsection{Romtoken for videostrømmer (S1.2)}
\label{subsec:krav-s12}

Hvert videorom skal tildeles en sikkerhetstoken slik at ikke uvedkommende 
kan få tilgang til strømmen.

\textbf{Rasjonale:} Når AMK-operatør sender SMS-lenke til innringer/pasient, 
må lenken inneholde en token som sikrer at kun mottaker kan delta i 
videorommet. Dette forhindrer at uvedkommende kan gjette rom-ID og få 
tilgang til sensitiv video.

\textbf{Implementasjon:} AntMedia bruker JWT-baserte viewer tokens, mens 
NHN Video bruker Pexip Guest PINs. VideoAPI abstraherer disse forskjellene 
og returnerer en uniform tilkoblingslenke uavhengig av backend.